\documentclass[11pt]{article}

\usepackage[a4paper,margin=1in]{geometry}
\usepackage{amsmath,amssymb,amsfonts}
\usepackage{enumitem}
\usepackage{hyperref}

\begin{document}

\begin{center}
{\Large \textbf{CSE400 – Fundamentals of Probability in Computing}}\\[4pt]
{\large Lecture 5: Bayes’ Theorem, Random Variables, and Probability Mass Function}\\[4pt]
Date: January 20, 2026
\end{center}

\section*{Definitions and Notation}

\subsection*{Events}
Let $A$ and $B$ be events.

\subsection*{Event Decomposition}
\[
A = AB \cup AB^{c}
\]

\subsection*{Conditional Probability}
\[
\Pr(A \mid B), \quad \Pr(A \mid B^{c})
\]

\subsection*{Random Variable}
A real-valued function defined on the sample space. Values are determined by outcomes of an experiment. Probabilities are assigned to possible values of random variables.

\subsection*{Discrete Random Variable}
A random variable that can take on at most a countable number of possible values.

\subsection*{Range of a Discrete Random Variable}
Let $X$ be a discrete random variable with range
\[
R_X = \{x_1, x_2, x_3, \dots\}
\]
(finite or countably infinite).

\subsection*{Probability Mass Function (PMF)}
The function
\[
p(x_k) = \Pr(X = x_k)
\]
is called the \textbf{Probability Mass Function (PMF)} of $X$.

\section*{Assumptions}
\begin{itemize}[leftmargin=1.5em]
\item Events $AB$ and $AB^{c}$ are mutually exclusive.
\item Axiom 3 of probability applies for mutually exclusive events.
\item Probabilities of all possible values of a discrete random variable sum to 1.
\item Coins and cards mentioned in examples are fair unless stated otherwise.
\end{itemize}

\section*{Theorems / Results}

\subsection*{Weighted Average of Conditional Probabilities}
For events $A$ and $B$,
\[
\Pr(A) = \Pr(AB) + \Pr(AB^{c})
\]
\[
\Pr(A) = \Pr(A \mid B)\Pr(B) + \Pr(A \mid B^{c})[1 - \Pr(B)]
\]

\subsection*{Law of Total Probability (Formula 3.4)}
If $\{B_i\}$ forms a partition of the sample space, then
\[
\Pr(A) = \sum_i \Pr(AB_i)
\]

\subsection*{Bayes Formula (Proposition 3.1)}
Using
\[
\Pr(AB_i) = \Pr(B_i \mid A)\Pr(A)
\]
we get
\[
\Pr(B_i \mid A) =
\frac{\Pr(A \mid B_i)\Pr(B_i)}{\sum_j \Pr(A \mid B_j)\Pr(B_j)}
\]

\section*{Proofs / Derivations}

\subsection*{Weighted Average of Conditional Probabilities}
Since
\[
A = AB \cup AB^{c}
\]
and $AB$ and $AB^{c}$ are mutually exclusive,
\[
\Pr(A) = \Pr(AB) + \Pr(AB^{c})
\]

Using conditional probability definitions,
\[
\Pr(AB) = \Pr(A \mid B)\Pr(B)
\]
\[
\Pr(AB^{c}) = \Pr(A \mid B^{c})\Pr(B^{c})
\]

Thus,
\[
\Pr(A) = \Pr(A \mid B)\Pr(B) + \Pr(A \mid B^{c})[1 - \Pr(B)]
\]

\subsection*{Bayes Formula}
Substitution into the law of total probability yields the Bayes Formula. Further derivation steps are omitted.

\section*{Worked Examples}

\subsection*{Example 3.1 (Part 1)}
\[
\Pr(\text{accident} \mid \text{accident prone}) = 0.4
\]
\[
\Pr(\text{accident} \mid \text{not accident prone}) = 0.2
\]
\[
\Pr(\text{accident prone}) = 0.3
\]

Let $A_1$: accident within a year, $A$: accident prone.

\[
\Pr(A_1) = (0.4)(0.3) + (0.2)(0.7) = 0.26
\]

\subsection*{Example 3.1 (Part 2)}
\[
\Pr(A \mid A_1) = \frac{(0.3)(0.4)}{0.26} = \frac{6}{13}
\]

\subsection*{Example 3.2 (Cards)}
\[
\Pr(RB \mid R) =
\frac{(1/2)(1/3)}{(1)(1/3) + (1/2)(1/3) + (0)(1/3)} = \frac{1}{3}
\]

\section*{Random Variable Example}

Experiment: Tossing 3 fair coins.

\[
\Pr(Y=0)=\frac{1}{8},\quad
\Pr(Y=1)=\frac{3}{8},\quad
\Pr(Y=2)=\frac{3}{8},\quad
\Pr(Y=3)=\frac{1}{8}
\]

\[
\sum_{y=0}^{3}\Pr(Y=y)=1
\]

\section*{PMF Example}

\[
p(i)=c\lambda^i,\quad i=0,1,2,\dots
\]

Find $\Pr(X=0)$ and $\Pr(X>2)$. Steps are stated in the lecture material; derivations are omitted.

\section*{Conclusion}

\begin{itemize}[leftmargin=1.5em]
\item Probability of an event can be expressed as a weighted average of conditional probabilities.
\item Law of Total Probability and Bayes Formula relate conditional and marginal probabilities.
\item Random variables map outcomes to real values.
\item Discrete random variables are characterized by probability mass functions.
\end{itemize}

\end{document}